% !PRINTED 238
\weberpage{238}
%<<CONT
von denen die beiden ersten
\[
1)\quad D_n' = \begin{pmatrix} e^{\frac{\pi i h}{2n}}, & 0 \\ 0, & (-1)^h e^{-\frac{\pi i h}{2n}} \end{pmatrix},\quad
\begin{pmatrix} 0, & -e^{\frac{\pi i h}{2n}} \\ (-1)^{h+n} e^{-\frac{\pi i h}{2n}}, & 0 \end{pmatrix}
\]
\[
h = 0, 1, \ldots, 2n-1,
\]
\[
2)\quad D_n'' = \begin{pmatrix} e^{\frac{\pi i h}{n}}, & 0 \\ 0, & \pm\, e^{-\frac{\pi i h}{n}} \end{pmatrix},\quad
\begin{pmatrix} 0, & e^{\frac{\pi i (2h+n)}{2n}} \\ \pm\, e^{-\frac{\pi i (2h+n)}{2n}}, & 0 \end{pmatrix}
\]
\[
h = 0, 1, \ldots, n-1
\]
sind, während die dritte Gruppe (9)
\[
\begin{pmatrix} e^{\frac{\pi i h}{n}}, & 0 \\ 0, & e^{-\frac{\pi i h}{n}} \end{pmatrix},\quad
\begin{pmatrix} 0, & e^{\frac{\pi i h}{n}} \\ e^{-\frac{\pi i h}{n}}, & 0 \end{pmatrix},
\]
\[
\begin{pmatrix} 0, & e^{\frac{\pi i}{2n}(2h+n)} \\ -e^{-\frac{\pi i}{2n}(2h+n)}, & 0 \end{pmatrix},\quad
\begin{pmatrix} e^{\frac{\pi i}{2n}(2h+n)}, & 0 \\ 0, & -e^{-\frac{\pi i}{2n}(2h+n)} \end{pmatrix}
\]
\[
h = 0, 1, \ldots, n-1
\]
ergiebt, was bei geradem $n$ mit $D_n''$, bei ungeradem $n$ mit $D_n'$ übereinstimmt.

Der Fall $n=2$, wo $c = \begin{pmatrix} i, & 0 \\ 0, & -i \end{pmatrix}$ ist, bildet nur eine scheinbare Ausnahme, weil in diesem Falle
\weq{10}{\varphi^{-1} c \varphi = d \quad \text{oder} \quad \varphi^{-1} c \varphi = c d}
sein kann. Verfolgt man diese Annahmen durch einfache Rechnung, so ergeben sich noch zwei Gruppen:
\weq{11}{
\begin{aligned}
&1,\ c,\ d,\ cd,\ \varphi_1,\ \varphi_1 c,\ \varphi_1 d,\ \varphi_1 c d\\
&1,\ c,\ d,\ cd,\ \varphi_2,\ \varphi_2 c,\ \varphi_2 d,\ \varphi_2 c d,
\end{aligned}
}
worin
\[
\varphi_1 = \begin{pmatrix} \dfrac{1}{\sqrt2}, & \dfrac{1}{\sqrt2} \\ \dfrac{1}{\sqrt2}, & \dfrac{-1}{\sqrt2} \end{pmatrix},\quad
\varphi_2 = \begin{pmatrix} \dfrac{1}{\sqrt2}, & -\dfrac{i}{\sqrt2} \\ \dfrac{i}{\sqrt2}, & -\dfrac{1}{\sqrt2} \end{pmatrix}.
\]
